% s01b_intro (abridged version). Introduction, part 2: inputs, what is not claimed, strategy, plan.

%-------------------------------------------------------------------------------
\subsection{Results from the literature used in the proofs}
\label{sec:int-inputs}

The proofs of \cref{thm:int-frame,thm:int-regularity,cor:int-gap,cor:int-gapkh} use the following published results. Their proofs are not reproduced here or in \textup{[L]}.
\begin{description}
\item[Mirror theorem at generic weights] \cite[Theorem~31]{CCIT15}, with Definition~28, eq.~(19), eq.~(4) and Remark~3 there, applied to $Z$ in \cref{thm:gen-mirror}. The genus-zero differential equation at generic weights is not a further input as a whole: it is derived in \cref{prop:gen-qde} from \cite[Lemma~1]{Pan98} and the boundary structure.
\item[Classification of semisimple theories] \cite[Theorem~(2), Proposition~8.7(ii), 2.17--2.19, (2.15)]{Tel12}, which rest on Harer stability, as stated in \cite[Propositions~(2.1), (2.2)]{Loo96} with \cite{Har85,Iva90}, and on \cite[Theorem~1.1.1]{MW07}.
\item[The group action] \cite[Definitions~2.2, 2.5, 2.13, Proposition~2.14, Section~3.2 and (5)]{PPZ15}.
\item[Localisation and the virtual class] \cite[formula~(1), Appendices~A, C]{GP99}; the fundamental-class and splitting axioms of the virtual class, with \cite[Appendix~1 (iv)--(vi)]{CG07} for the twist.
\item[Hodge classes] Mumford's formula in the form of \cite[Theorem (Mumford), Section~1.1]{FP00}; the relation $c(\Hdg)c(\Hdg^\vee)=1$ of \cite[(5.4), (5.5)]{Mum83}; and \cite[Theorem~4]{FP00}.
\item[Special functions] \cite[Eq.~15.6.7]{DLMF}; Barnes's second lemma \cite{Bar10}; Gauss's values of the digamma function at $1/3$, $2/3$.
\item[Mari\~no--Vafa formula] \cite[eq.~(1)--(3)]{LLZ03}; \cite[Theorem~1]{OP04} is a second published proof. With it, the textbook facts: the string and dilaton equations with classes pulled back along forgetful maps, the genus-zero multinomial formula for $\psi$-integrals, the Frobenius formula and Frobenius' transposition identity, the values $\int_{\Mbar_{1,1}}\psi=\int_{\Mbar_{1,1}}\lambda_1=1/24$, the Murnaghan--Nakayama rule, the hook-length form of the principal specialisation of Schur functions, and Lagrange inversion.
\end{description}
Numbers quoted from \cite{CCIT15} and \cite{Tel12} are those of the arXiv versions 5 and 3. The results enter as follows. \Cref{thm:int-regularity} uses only the Mari\~no--Vafa formula and its textbook facts; the evaluation of $\ZZ[g,0]$ uses that formula only at $\mu=1^a$. The mirror theorem enters only through \cref{thm:int-frame}, and through it \cref{cor:int-gap,cor:int-gapkh}. The frame change, \cref{thm:fra-change}, does not use it. In \cref{lem:fra-tcf} the logarithmic form of $\wcf$, which gives $\Fcf=\theta\log\ccf$ in \cref{thm:int-frame}, uses \cite[Eq.~15.6.7]{DLMF} and Gauss's values, and Barnes's second lemma is needed only for the constant in $\tcf=\tc$.

We use \cite[Theorem~31]{CCIT15} at generic weights because \cref{sec:gen} cannot be moved to $\TF$: the quantum product at $\cQ=0$ is not semisimple there. The theorem of \cite{Coa09} gives the mirror statement on $\TF$. No topological recursion on the mirror curve of local $\bbP^2$ is used, and neither is remodelling. The only relation in the paper that is called a topological recursion is the genus-zero psi-class relation of \cref{prop:gen-qde}(i). The proof of \cite[Theorem~31]{CCIT15} does not use such input either; see \Lg{Remark~4.12}. We also do not use the conifold leading term, the pole bound or the graph sum for $\beta$ of \cite{BS}; \Lg{Remarks~10.13, 10.14} compare them with the present results. No proof step is a finite-order computation; evidence of that kind is collected in \Lg{Appendix~A}, see \cref{sec:fin}, and no proof uses it.

%-------------------------------------------------------------------------------
\subsection{What is not claimed}
\label{sec:int-notclaimed}

We claim nothing off the fibrewise torus, and in particular nothing about polar parts at general Calabi--Yau tori. We do not determine the value of the regular part of $\GWcf_g$. We do not show that $\sF_{k,h}$ lies in $\bbQ[X^{\pm1}][F]$: we prove only that $\sF_g^{\TF}$ lies in $\bbK[y^{\pm1},\xi^{-1}][F]$, with coefficients of $\sF_{k,h}$ in $\bbQ[X^{\pm1},(1-X)^{-1}]$, and the stronger form remains open for all $g$; it is independent of \cref{cor:int-gap}. We do not address the orbifold regularity, the holomorphic anomaly equations or the uniqueness statement in refined direct integration. The relation of $\Cf$ to a change of polarisation of the Givental quantisation is not addressed either. The cited theorems of \cref{sec:int-inputs} stay cited.

%-------------------------------------------------------------------------------
\subsection{Strategy and plan}
\label{sec:int-strategy}

The proof has two independent parts: an identification of $\GWcf_g$ with the explicit sum $\Gsum_g$ and an analysis of $\Gsum_g$ at the conifold point. The sections follow those of \textup{[L]}.

\emph{The frame formula} (\cref{sec:gen,sec:spe,sec:fra}). At generic weights on the maximal torus the theory is semisimple at $\cQ=0$, and \cite{Tel12} gives $\Omega=\Rm.\TFT$. The rows of $\Rm$ come from the mirror theorem and from the genus-zero differential equation, flatness of $\Rm$ is derived directly, and the constants are fixed by the point $\cQ=0$ (\cref{thm:gen-main}). The sum over colourings at generic weights is a multi-trace over the local ring $\ringC=\locO[Y]/(\Lhat(Y))$, with $\locO$ the localisation at the prime of $\TF$. So the graph sum is a polynomial in an indeterminate $\sfF$ with coefficients in $\locO$, and it passes to $\TF$ by a ring map, with no limit in the weights (\cref{thm:spe-main}); the theorem of \cite{Tel12} is not applied on $\TF$. On $\TF$ the factor $\hb=\Theta(-z/E)\Theta(z/\epsilon_1)\Theta(z/\epsilon_2)$ is the same at all colours, and by Mumford's formula it can be moved into the vertices as the class $\Hg_{\bw}$ (\cref{thm:fra-change}). This is an identity for the total, polynomial in $\sfF$, and the substitution $\sfF\mapsto\Fcf$ gives \cref{thm:int-frame}.

\emph{Regularity} (\cref{sec:row,sec:lea,sec:van,sec:val,sec:reg}). At the singular colour, $\Delta_0=3\epsilon_1\epsilon_2E^3/(\tc\rho)^2$ with $\rho$ an analytic unit, and to every order in $\tc$ the singular row and the factors built from it are finite combinations of the leaves $\lf_j(\psi)=1/(E-j\psi)$, with $j\leq k+1$ at order $\tc^k$ (\cref{thm:lea-main}). This is algebra on the closed form of the rows; it uses the identity \eqref{eq:row-zinf} and the symplecticity of the rows, \cref{prop:row-sympl}. By homogeneity under the dilaton equation, a colour-$0$ vertex of genus $h\geq 2$ then contributes $\tc^{2-2h+k}$ times finite combinations of the numbers $\Xiz_h(\msU)$ with $\dg(\msU)=\sum_{j\in\msU}(j-1)\leq k$; these are Hodge integrals with the class $\Hg_{\bw}$, and they involve neither the curve nor the periods of local $\bbP^2$. They vanish for every non-empty multiset $\msU$ of positive integers with $\dg(\msU)\leq 2h-3$ (\cref{thm:van-main}, \cref{cor:van-parts}, \cref{lem:reg-A}): the pair of Mari\~no--Vafa series related by the shift $p_k\mapsto\re^{k\tau x}p_k$ satisfies one explicit bilinear equation for each $k\geq 2$, and the lowest coefficient of that equation at the dilaton point gives an all-genus relation among these numbers. A count of orders then gives regularity of every graph with an edge, colouring by colouring (\cref{thm:reg-main}). In the graph with one vertex the empty multiset survives, with the value $\ZZ[g,0]$, which is computed from the Mari\~no--Vafa formula at $\mu=1^a$ and the first equation of the hierarchy (\cref{thm:val-main}).

The table of dependencies between the results and the literature is in \Lg{Section~1.5}. In particular, \cref{sec:lea,sec:van,sec:val} and \cref{thm:reg-main,cor:reg-sv} can be read independently of \cref{sec:gen,sec:spe,sec:fra}. \Cref{sec:fin} records the counterpart for the resolved conifold, which is conditional on a statement that is not proved in \textup{[L]}; it summarises the finite-order evidence, and it lists what the results are conditional on, including the checks of citations that \textup{[L]} marks as outstanding.

\subsection*{Acknowledgements}

This work was supported by the UKRI grant UKRI2394 and by the Leverhulme Trust grant LIP-2023-007.
